% kind: home
% title: 潘小工
% nav_title: 00-关于本站
% author: 潘小工
% date: 2026-03-06 00:00:00
% summary: 山东大学密码科学与技术专业本科生，热爱数学的密码学生。

\begin{center}
\begin{figure}
\includegraphics[width=0.2\textwidth]{https://github.com/panxiaogong.png}
\end{figure}

China/Qingdao \quad \href{https://github.com/panxiaogong}{GitHub}
\end{center}

\section{关于笔者}

你好，我是潘小工，是一个热爱数学的密码学生，来自山东大学网络空间安全学院，主修密码科学与技术，辅修金融数学。

如果你也在同一条学习路上，欢迎交流。

\section{关于本站}

我在上学院的数学课的时候，总是会想，能不能用最通俗易懂的语言来把这门课讲清楚，能不能从其他的优秀教材中获得更好的理解。结合这两个方面，我制作了这个网站，不仅仅是为了实现自己的想法，也希望能帮助到其他同学。

本站目前主要涵盖以下课程内容：

除数学课程外，后续也会逐步整理计算机相关课程的笔记。

这些笔记均使用 LaTeX 编写，你可以在 GitHub 仓库中找到对应的源码和 PDF 文件。

\section{打开方式}

不同的课程在侧边栏中会有不同的编号：

\begin{itemize}
\item 课内数学课程对应的编号为 \textbf{AXX}
\item 课内计算机课程对应的编号为 \textbf{CXX}
\item 课外数学书籍对应的编号为 \textbf{EXX}
\item 与数学工具相关的内容为 \textbf{TXX}
\end{itemize}

点击侧边栏中的选项即可展开对应课程的子目录，点击右侧区域中的具体条目即可打开正文。右侧目录可在当前页面内快速定位小节。

\section{网站来由}

我在大一下学期的时候，学离散数学学得非常开心，就想能不能整理自己的离散笔记，看看能不能整理得清楚明白，并顺手开源到 GitHub 的仓库里。结果非常成功，有好多同学都来看我的笔记，并且在期末周用我的笔记复习，我收获了很多好评和催更，这也成了我做这件事的一个小小转折点。

后来我发现 GitHub 的 Markdown 语法和主流语法有些不太兼容，尤其是数学公式。我当时用的 Markdown 编辑器是 Typora，当我把编辑好的 Markdown 格式复制到 GitHub 页面里面的时候，有些语法不得不重复修改，搞得我很恼火。干脆不用 GitHub 的页面了，自己做了一个博客来放这些笔记。

但这个博客仍然是基于 GitHub 的，而且使用了别人开发好的开源工具 Hexo。使用了几个月下来发现并不好用，一是因为 GitHub 的服务器在国外，访问速度不理想；二是因为 Hexo 的风格实在不合我的期望。所以我干脆从零开始，在 AI 辅助下做一整个框架——AI 实现框架，我输出内容，很快得到了我满意的一版，就是大家现在看到的这一版。

但是，这还不够。我只允许 AI 做框架，不允许 AI 来生成内容。AI 的学习路径与人类不同，让 AI 生成讲解文档我觉得是一种不负责任的行为，所以所有内容都是我自己来写。想要通俗易懂地讲清楚真的太难了，往往花费了难以想象的时间只能得到很少的产出，但我觉得这是值得的。

以上是这个网站的历史，下面我来说点我自己想说的。

我不是一个很有天赋的人，在过去五年的努力里，我知道自己已经来到了自己的上限。尤其是在我写这些文档的时候，我不知道该怎么形容这种感觉——有点穷途末路的无奈。但是我知道，屏幕前的各位还远没有到达自己的上限，还远没有意识到自己的潜力，而我愿意相信大家的能力。

沉睡的狮子需要时间来觉醒，我去完成最后的一点铺垫。